% 本文件由 LaTeX 排版 Agent 根据有效 translation.json 自动生成。
% author=Liturgies; work_id=0719; title=圣阿代乌斯与圣玛里斯礼仪

\chapter{圣阿代乌斯与圣玛里斯礼仪}

% structure: p0001 source=h1 target=omitted reason=page-title-already-rendered-by-parent

\Emph{由圣阿代乌斯与圣玛里斯，东方诸导师所作。}\par

一、首先：光荣归于至高天主，等等。我们的天父。\par

\Emph{祈祷词。}\par

主，我们的天主，求祢因祢的仁慈坚固我们的软弱，使我们能举行为更新及拯救我们堕落本性而赐予的圣奥迹，藉着祢爱子、万有之主的仁慈。\par

\Emph{平日。}\par

万有之主，愿祢永受赞颂的天主圣三那可敬而荣耀的圣名，在天地间受到朝拜、荣耀、赞颂、颂扬、高举及祝福。\par

\BibleBook{}{圣咏}第十五篇：\BibleRef{咏 15:1} \BibleQuote{主，谁能寄居在祢的帐幕里？}全文及其圣咏曲，关联圣事奥迹。\par

\Emph{（高声。）}\par

谁能欢呼？等等。\par

\Emph{祈祷词。}\par

二、主，在祢威严的光辉宝座前，在祢荣耀的崇高至伟的宝座前，在祢爱情的威严可畏之座前，在祢旨意所立的赎罪祭台前，在祢牧场的区域中，偕同千万赞美祢的革鲁宾与亿万颂扬祢的色辣芬，我们前来朝拜、感谢并永远荣耀祢，万有之主。\par

\Emph{纪念日及星期五。}\par

万有之主，祢伟大而圣洁、光辉且受颂扬的圣名，祢荣耀圣三的受颂扬而不可理解的圣名，以及祢对我们民族的恩慈，我们应当时时刻刻赞美、朝拜并荣耀祢。\par

\Emph{在圣所前的应答咏，如上。}\par

那位吩咐者，等等。\par

向司铎，等等。\par

\Emph{祈祷词。}\par

主，我们的天主，祢爱之甘饴的馨香如何在我们内呼吸；我们的灵魂因认识祢的真理而得光照。求使我们堪当从祢的圣天领受祢爱子的显现；我们将在那里感谢祢，并在现时于祢的教会中不停地荣耀祢，教会充满各种助佑与祝福，因祢是主，是父，万有的创造者。\par

\Emph{三、献香祈祷词。}\par

我们将向祢荣耀的天主圣三，圣父、圣子、圣神，重复赞歌。\par

\Emph{守斋日。}\par

并且因着，等等。\par

\Emph{在纪念圣人时。}\par

主，祢确实是使我们身体复起者，是我们灵魂的善牧，是我们生命的安全守护者；我们应当时常感谢祢、朝拜并荣耀祢，万有之主。\par

\Emph{在诵读时。}\par

祢是圣洁的，堪受赞颂，全能，永生，居住在圣者之中，祢的旨意安息于他们内。主，求祢垂顾我们，求祢怜悯我们，怜恤我们，因为祢是我们一切处境中的援助，万有之主。\par

\Emph{四、在宗徒书信时。}\par

主，我们的天主，求祢光照我们默想的活动，使我们听到并理解祢赋予生命与神圣命令的甘美之言；并藉祢的恩宠与仁慈，赐我们从中获得爱情、望德及合宜于灵魂与身体的救恩的确证；我们将不停地永远向祢咏唱永恒的荣耀，万有之主。\par

\Emph{守斋日。}\par

向着祢，智慧的治理者，等等。\par

\Emph{五、下台，向福音致礼，在祭台前诵念此祈祷词。}\par

祢，祢圣父的卓越后裔，祢圣父位格的肖像，曾在人性身体中显现，并藉祢预报之光向我们升起，我们感谢祢、朝拜祢，等等。\par

\Emph{并在宣告后：——}\par

主，全能的天主，我们恳求并祈求祢，以祢的恩宠成全我们，并藉我们的手倾注祢的恩赐，即祢神性的怜悯与慈悲。愿这些为我们成为祢子民过犯的赎罪，并为祢牧场全体羊群罪过的赦免，藉祢的恩宠与仁慈，良善的世人爱者，万有之主。\par

\Emph{六、执事们说：——}\par

低头。\par

\Emph{司铎在至圣所念此秘密祈祷词：——}\par

全能的主天主，圣而公教会是祢的，因祢藉祢基督的伟大苦难赎买了祢牧场的羊群；并从圣神的恩宠（祂确实与祢荣耀的神性同体）赐予了真正司铎圣秩的等级；并藉祢的仁慈，主，祢屈尊使我们这些软弱者在祢圣教会的大身体中成为属灵的肢体，为能向信友灵魂施行属灵的援助。如今，主，求祢以祢的恩宠成全我们，并藉我们的手倾注祢的恩赐；愿祢的怜悯与祢神性的仁慈临于我们，并临于祢为自己所拣选的子民。\par

\Emph{（高声。）}\par

主，求祢藉祢的仁慈恩赐我们，使我们众人能一同并平等地在我们生命的每一天令祢的神性喜悦，并堪当获得祢恩宠的助佑，时时向祢献上赞美、尊荣、感谢及朝拜，主。\par

\Emph{七、然后执事们登上祭台，说：——}\par

未领受圣洗者，等等。\par

\Emph{然后司铎开始奥迹的应答咏，圣器管理员与执事将圣盘与圣爵置于祭台上。司铎交叉双手，说：——}\par

我们向祢荣耀的天主圣三献上赞美，时时直到永远。\par

\Emph{然后继续：——}\par

愿那为我们的救恩而被献，并命令我们纪念祂死亡与复活的基督，亲自从我们软弱的手中接纳此祭献，藉祂的恩宠与仁慈直到永远。阿们。\par

\Emph{然后继续：——}\par

那著名的、圣洁而赋予生命的奥迹已安放在大能之主的祭台上，直到祂的来临，直到永远。阿们。\par

赞美，等等。\par

祢的纪念，等等。\par

我们的天父，等等。\par

圣父的宗徒们，等等。\par

在圣祭台上，等等。\par

已安眠者，等等。\par

玛窦、玛尔谷、路加，等等。\par

\Emph{信经。}\par

\Emph{八、司铎走近举行祭献，三次俯伏于祭台前，吻其中间，然后吻祭台右角与左角；并向福音侧俯首，说：——}\par

主，求祢降福，等等。\par

请为我祈祷，我的父老们、弟兄们及师长们，求天主赐我能耐与权柄，举行我现已走近的这项礼仪，并愿此祭献从我软弱的手中蒙接纳，为我自己、为你们、并为整个圣而公教会，藉祂的恩宠与仁慈直到永远。阿们。\par

\Emph{他们回应：——}\par

愿基督俯听你们的祈祷，悦纳你们的祭献，收纳你们的奉献，尊荣你们的司祭职，并因你们的转求，赐予我们过犯的赦免和罪恶的宽恕，藉着祂的恩宠和仁慈，至于永远。\par

\Emph{随即他在另一边鞠躬，说出同样的话；他们也同样回应：然后他向祭台鞠躬，说：——}\par

天主，万有的主，愿祂藉着祂的恩宠和仁慈常与我们同在。阿们。\par

\Emph{并向左边的执事鞠躬，说：——}\par

天主，万有的主，坚定你的言语，赐你平安，并从我手中为我自己、为你、为圣而公教会全体、并为整个世界，收纳这祭品，藉着祂的恩宠和仁慈，至于永远。\par

\Emph{他跪在祭台前，秘密地说：——}\par

九、我们的主、天主，不要看我们的众多罪过，勿因我们过犯的凶恶而转离祢的尊威；但藉祢不可言喻的恩宠，圣化祢的这祭献，并藉它赐予力量和效能，使祢能忘却我们许多的罪，当祢在时期终结时，以从我们中取的人性显现时，能仁慈待我们；使我们能寻得祢面前的恩宠和仁慈，并堪当与天上的会众一起赞美祢。\par

\Emph{他起身，秘密地说这祷词：——}\par

我们感谢祢，我们的主、天主，为祢对我们丰富的恩宠之财宝：\par

\Emph{他继续：——}\par

我们这些因祢大慈大悲而得以犯罪堕落的人，竟蒙祢使我们堪当举行祢基督体血的圣奥迹。我们向祢祈求援助，以坚固我们的灵魂，使我们能以完美的爱和真实的信德，施与祢赐予我们的恩赐。\par

\Emph{感恩经。}\par

我们将把赞美、荣耀、感恩和朝拜归于祢，从今时到永远，及世之世。\par

\Emph{他划十字圣号，他们回应：——}\par

阿们。\par

\Emph{十、他继续：——}\par

愿平安与你们同在。\par

\Emph{他们回应：——}\par

也与你的心灵同在。\par

\Emph{他们彼此行平安礼，并说：——}\par

为众人：\par

\Emph{执事说：——}\par

让我们感谢、恳求并祈求。\par

\Emph{司铎秘密地说这祷词：——}\par

主，大能的天主，求祢因祢的仁慈和祢恩宠的助佑，扶助我的软弱；并使我堪当在祢面前奉献这祭品，作为众人的共同援助，并赞美祢的圣三，圣父、圣子、圣神。\par

\Emph{另一祷词。}\par

我们的主、天主，约束我们的思念，使它们不游荡于世间的虚妄中。主、我们的天主，求祢使我，虽不配，能与祢爱之情结合。愿光荣归于祢，基督。\par

主，请上升至祢荣光的居所；在我内播下谦逊的善种；并以祢恩宠的翅膀荫庇我，因祢的仁慈。主，祢若细察罪过，谁能站立？因为祢有仁慈。\par

\Emph{[司铎秘密地说以下祷词：——}\par

我们主耶稣基督的母亲，请为我恳求那由你而生的独生子，宽恕我的过犯和罪恶，并从我软弱有罪的双手，接受这卑微的祭献，这祭献由我的软弱在这祭台上奉献，藉着你为我的转求，圣母。]\par

\Emph{十一、当执事说\Quote{以警醒和谨慎}等时，司铎立即起身，揭开圣事，除去覆盖的帷幔：他祝福乳香，并高声诵念感恩经：——}\par

愿我们主耶稣基督的恩宠，天父的爱，和圣神的共融，与你们众人同在，从今时到永远。\par

\Emph{他划十字圣号，他们回应：——}\par

阿们。\par

\Emph{司铎继续：——}\par

请高举你们的心。\par

\Emph{他们回应：——}\par

我们已高举向祢，亚巴郎、依撒格、以色列的天主，荣耀的君王。\par

\Emph{司铎。}\par

祭品已献给天主，万有的主。\par

\Emph{他们回应：——}\par

这是应当的、正义的。\par

\Emph{执事。}\par

愿平安与你们同在。\par

\Emph{司铎放乳香，并说这祷词：——}\par

主，上主，求祢在祢面前赐我坦然的面容，使我能因从祢而来的自信，以无罪、无苦毒的良心，举行这可畏的神圣祭献。主，求祢在我们内播下彼此之间及对所有人的情感、和平与和睦。\par

\Emph{他站立，秘密地说：——}\par

堪受一切口舌的荣耀、一切唇舌的感恩、一切受造物的朝拜和颂扬的，是圣父、圣子、圣神那可敬可颂的圣名，祂以恩宠创造了世界，以仁慈创造了世上的居民，以怜悯拯救了人类，向必朽者显示了极大的恩惠。主，祢的威严，千万层天上的神体，与万万千千的圣天使、众神之军、火与神之臣仆，都祝福朝拜；与圣革鲁宾和神性塞拉芬，他们圣化并欢庆祢的圣名，呼喊赞美，不停彼此呼喊。\par

\Emph{他们高声说：——}\par

\BibleQuote{圣、圣、圣，主、全能的天主；祂的光荣充满天地。}\BibleRef{依 6:3}\par

\Emph{司铎秘密地说：——}\par

圣、圣、圣，祢是主、全能的天主；天地充满祂的光荣和祂本性的实质，正如它们以祂光辉的荣耀为荣；正如经上记载：\Quote{天地都充满了我，全能的主说。}\par

祢是圣的，天主我们的父，真正独一的那位，天上地下的一切家族都由祂得名。祢是圣的，永恒的圣子，万物藉着祂造成。祢是圣的，圣的、永恒的圣神，万物藉着祂被圣化。\par

\BibleQuote{祸哉，我，祸哉，我，我已惊惶失色，因为我是个唇舌不洁的人，住在唇舌不洁的人民中，我的眼睛看见了君王，全能的主。}\BibleRef{依 6:5}\newline{} \BibleQuote{今日这地方多么可畏！这实非他处，乃是天主的殿，天门；因为主，祢已面对面被看见。}\BibleRef{创 28:17}\par

如今，我祈祷，愿祢的恩宠与我们同在，主；洗涤我们的不洁，圣化我们的口唇；将我们卑微的呼声与塞拉芬和总领天使的圣化联合。愿光荣归于祢的仁慈，因为祢使地上的与天上的结合。\par

\Emph{他继续前行，俯身暗自诵念此祷文：—}\par

十二、我们与那诸天的神力一同感谢祢，我们这些微不足道、软弱无能的仆人，因为祢赐予了我们无可报答的浩大恩宠。祢确实取了我们的人性，好藉祢的天主性赐予我们生命；祢高举了我们卑微的处境；祢复兴了我们颓丧的状态；祢唤醒我们必朽的生命；祢洗净了我们的罪；祢抹去了我们罪过的债；祢光照了我们的理智，并谴责了我们的仇敌——上主，我们的天主；祢使我们软弱的本性得胜。\par

\Emph{此处为建立圣体之言，随后：—}\par

请藉祢倾注恩宠的柔情，宽仁之主，赦免我们的过犯与罪愆；在审判中抹去我的过犯。因祢赐予我们的一切助佑与恩惠，我们将向祢献上赞美、尊荣、感恩与朝拜，从今时直到永远。\par

\Emph{司铎在圣事上划十字圣号。会众回应。}\par

阿们。\par

\Emph{执事：}\par

请你们心中思念。请与我们一同祈求平安。\par

\Emph{司铎躬身低声诵念此祷文：—}\par

上主，全能的天主，请接纳这祭品，为整个神圣天主教会，并为所有蒙祢悦纳的虔诚祖先，以及为所有先知与宗徒，并为所有殉道者与宣信者，并为所有哀恸者、困苦者、患病者，并为所有遭患难与试探者，并为所有软弱与受压迫者，并为所有已离世的亡者；然后为所有从我们软弱中祈求祷告的人，并为我这卑微软弱的罪人。上主，我们的天主，请照祢的仁慈与丰厚的恩惠，垂顾祢的子民，并垂顾我这软弱之人，不要按我的罪过与愚妄，而是使他们因着信德领受这圣体，藉祢仁慈的恩宠，得以堪当罪过的赦免，直到永远。阿们。\par

\Emph{司铎暗自诵念此俯首祷文：—}\par

十三、上主，求祢藉祢丰厚的、不可言喻的仁慈，使这纪念与所有在祢面前祈求的虔诚祖先的纪念一同蒙悦纳，他们是在纪念祢的基督的圣体圣血时——我们按祢所教导的，在祢纯洁神圣的祭台上呈献给祢——求祢赐予我们在这生命的日子里安息。\par

\Emph{他继续大祭献：—}\par

上主，我们的天主，求祢赐予我们在这生命的日子里安息与平安，使地上的万民都认识祢，唯独祢是真天主圣父，并认识祢所派遣的主耶稣基督，祢的圣子，祢的爱子；祂亲自，我们的主与天主，来临并教导我们一切纯洁与圣洁。求祢纪念先知、宗徒、殉道者、宣信者、主教、圣师、司铎、执事，以及所有神圣天主教会中接受了生命之印——圣洗——的子女。我们也求祢，上主：\par

\Emph{他继续：—}\par

我们，祢卑微、软弱、无能的仆人，因祢的名聚集，现今侍立在祢面前，欣然接受了从祢而来的形式，赞美、光荣、颂扬、纪念并举行这伟大、可敬、神圣、天主的奥迹，即我们的主、救主耶稣基督的苦难、死亡、埋葬与复活。\par

愿祢的圣神降临，上主，并安息在祢仆人所献的这祭品之上，祝福并圣化它；愿这祭品成为我们，上主，过犯的赎罪、罪过的赦免、对从死者中复活的大希望，以及在天国里与所有蒙祂悦纳的人同享新生命。因祢对我们整个奇妙的救世计划，我们将向祢感恩，并在祢的教会——以祢基督的宝血所救赎的教会——中，张口欢颜，不住地光荣祢：\par

\Emph{正典。}\par

将赞美、尊荣、感恩与朝拜归于祢神圣、仁爱、赋予生命的圣名，从今时直到永远。\par

\Emph{司铎用十字圣号划圣体，会众回应：—}\par

阿们。\par

\Emph{司铎躬身亲吻祭台，先中间，后左右两边，并诵念此祷文：—}\par

\Quote{天主，求祢怜悯我}，直到\Quote{罪人将归向祢}；\Quote{上主，我向祢举目}，直到\Quote{上主，求祢怜悯我们，怜悯我们}。并伸出祢的手，愿祢的右手拯救我，上主；愿祢的仁慈常留于我，上主，直到永远，不要轻看祢双手的工程。\par

\Emph{然后他诵念此祷文：—}\par

十四、基督，天上者的平安，地下者的大安息，求祢使祢的安息与平安居于世界的四方，尤其是在祢神圣天主教会内；使司铎与政府享有和平；从地极止息战争，驱散好战的民族，使我们能享平静安宁生活的福分，在一切节制与敬畏天主中。求祢因祢的恩宠与仁慈，宽恕亡者的过犯与罪愆，直到永远。\par

\Emph{他对祭台周围的人说：—}\par

求主降福。求主降福。\par

\Emph{他加上乳香，熏自己，并说：—}\par

上主，我们的天主，求祢以祢爱情的甘饴，使我们的灵魂发出馨香，并以此洗净我罪的污点，赦免我的过犯与罪愆，无论我所知或所不知的。\par

\Emph{他再次双手取乳香，向圣体奉香；立即说：—}\par

祢恩宠的仁慈，我们的主与天主，使我们这些不配的人得以亲近这些著名、神圣、赋予生命、天主的奥迹。\par

\Emph{司铎重复这话一次二次，每次间隔将双手在胸前交叉成十字形。他亲吻祭台中央，双手取上面的祭饼；举目仰望说：—}\par

愿赞美归于祢的圣名，主耶稣基督，并朝拜归于祢的尊威，直到永远。阿们。\par

因为祂是自天降下的、活而赐予生命之粮，赐生命给整个世界，凡吃这粮的必不死；凡领受它的，必因它而得救，不见腐朽，并藉它而永远生活；祢是我们必死性的解毒剂，我们整个身躯的复活。\par

十五. * * *\par

十六. 赞美祢的圣名，上主。（同上。）\par

\Emph{司铎以十字形亲吻圣体，但嘴唇不接触，仅作亲吻状，并说：——}\par

光荣归于祢，上主；光荣归于祢，上主，因祢对我们不可言喻的恩赐，直到永远。\par

\Emph{随后他走近擘圣体，用双手完成，说：——}\par

上主，我们以真信德走近，以感恩并藉祢的仁慈，为赐生命者耶稣基督的体血划十字，因父、子及圣神之名。\par

\Emph{并呼号天主圣三之名，他把手中圣体擘为两半：左手中的放在圣盘上；右手中的用来在圣爵上划十字，说：——}\par

珍贵的宝血以我主耶稣基督的圣体划上十字。因父、子及圣神之名，直到永远。\par

\Emph{他们回应：——}\par

阿们。\par

\Emph{然后他把圣体浸入圣爵至一半，并用它在圣盘上的圣体上划十字，说：——}\par

圣体以我主耶稣基督的赎罪宝血划上十字。因父、子及圣神之名，直到永远。\par

\Emph{他们回应：——}\par

阿们。\par

\Emph{他将两半合在一起，说：——}\par

这些可敬、圣洁、赋予生命和神圣的奥迹，彼此间已被分开、圣化、完成、成全、联合、混合，在祢可敬的、光荣的、荣耀的天主圣三名下，圣父、圣子、圣神，使它们对我们，上主，成为我们过犯的补赎和罪的赦免；也为从死者中复活的伟大希望，以及在天国的新生命，为我们，也为基督我主的圣教会，无论在此处还是任何地方，现在、时常、直到永远。\par

十七. \Emph{其间他用右手拇指以十字形从下到上、从右到左在圣体上划十字，从而在浸过血的地方形成一道细裂痕。他将一部分圣体以十字形放入圣爵：下部朝向司铎，上部朝向圣爵，使裂痕处朝向圣爵。他鞠躬，然后起身说：——}\par

光荣归于祢，主耶稣基督，祢虽不配，仍藉祢的恩宠使我成为祢可敬、圣洁、赋予生命和神圣奥迹的仆役与中保：藉祢仁慈的恩宠，使我堪当获得我过犯的赦免和罪的宽恕。\par

\Emph{他在自己额上划十字圣号，也对周围站着的人同样做。}\par

\Emph{执事们上前，他在每人额上划十字，说：——}\par

基督悦纳你的职分：基督使你的面容发光：基督拯救你的生命：基督使你的青春成长。\par

\Emph{他们回应：——}\par

基督悦纳你的奉献。\par

\Emph{十八. 众人各自归位；司铎鞠躬后起身，以福音的声调说：——}\par

愿主耶稣基督的恩宠，天主圣父的爱，圣神的共融，与我们众人同在。\par

\Emph{司铎划十字，并举手过头，使手在空中，使民众参与咏唱：——}\par

\Emph{执事说：——}\par

我们众人怀着敬畏等。\par

\Emph{在这些话时：——}\par

祂已将祂的奥迹赐给了我们：\par

\Emph{司铎开始擘圣体，并说：——}\par

上主，求祢藉祢的仁慈，怜悯祢仆人们的罪过和愚妄，并藉祢的恩宠圣化我们的嘴唇，使它们能在祢的国度里，与祢的众圣徒一起，将光荣和赞美的果实献给祢的天主性。\par

\Emph{他提高声音说：——}\par

上主，我们的天主，求祢使我们堪当以纯洁的心、坦然的容颜，并藉祢仁慈赐予我们的依靠，常常无玷地侍立在祢面前；愿我们同心呼求祢，这样说：我们的天父等。\par

\Emph{民众说：——}\par

我们的天父等。\par

\Emph{司铎。}\par

上主，全能的天主，上主，我们的善神，充满仁慈，我们恳求祢，上主，我们的天主，并祈求祢良善的慈悯：不要让我们陷于诱惑，但救我们脱离那恶者及其党羽；因为国度、权能、力量、威能与统治，在天上及地上都属于祢，现在、时常。\par

\Emph{他划十字，他们回应：——}\par

阿们。\par

\Emph{十九. 他继续：——}\par

愿平安与你们同在。\par

\Emph{他们回应：——}\par

也与你的心灵同在。\par

\Emph{他继续：——}\par

圣物应归于圣者，臻于完美。\par

\Emph{他们说：——}\par

唯一圣父：唯一圣子：唯一圣神。光荣归于父、子及圣神，直到永远。阿们。\par

\Emph{执事。}\par

你们赞美。\par

\Emph{他们诵答唱咏。执事前来取圣爵时，他说：——}\par

让我们为我们之间的和平祈祷。\par

\Emph{司铎说：——}\par

圣神的恩宠与你们同在，与我们同在，并与领受祂的人同在。他把圣爵交给执事。\par

\Emph{执事说：——}\par

上主，请降福。\par

\Emph{司铎。}\par

愿我们生命之主耶稣基督恩宠的恩赐，在仁慈中与众人一同完成。\par

\Emph{他以十字架降福民众。其间诵答唱咏。}\par

弟兄们，请领受圣子的身体，教会呼吁，在祂国度的家中，以信德饮祂的圣爵。\par

\Emph{在庆节。}\par

上主，求祢坚固等。\par

\Emph{在主日。}\par

主耶稣基督等。\par

\Emph{平日。}\par

我们所领受的奥迹等。\par

\Emph{答唱咏结束后，执事说：——}\par

因此众人等。\par

\Emph{他们回应：——}\par

愿光荣归于祂，因祂不可言喻的恩赐。\par

\Emph{执事。}\par

让我们为我们之间的和平祈祷。\par

\Emph{司铎在祭台中央诵此祈祷：——}\par

XX. 主啊，在所有的日子、时刻和时辰，感谢、朝拜并赞美祢威严可畏的圣名，是合宜的、公义的，因为祢藉着祢的恩宠，主啊，使我们这些具有脆弱本性的有死之人，堪与天上的受造物一同圣化祢的圣名，并分享祢恩赐的奥迹，且因祢圣言的甘饴而喜乐。我们不断将荣耀与感恩的声音献于祢崇高的天主性，主啊。\par

\Emph{另一篇。}\par

基督，我们的天主、主、君王、救主和生命赋予者，藉着祂的恩宠使我们堪当领受祂的圣体以及祂宝贵而全圣的圣血。愿祂恩赐我们，使我们在言语、行为、思想和行动上都令祂喜悦，好使我们所领受的这个凭据，成为我们过犯的赦免、罪恶的宽恕、从死者中复活的伟大希望，以及在诸天之国内获得崭新而真实的生命，与所有在祂面前蒙悦纳的人在一起，藉着祂的恩宠和怜爱，直到永远。\par

\Emph{平日。}\par

赞美、荣耀、祝福与感恩，主啊，我们应将这一切归于祢荣耀的天主圣三，为祢所赐予我们的圣洁奥迹的恩赐，这些奥迹是祢赐给我们，为平息我们的过犯，万有的主啊。\par

\Emph{另一篇。}\par

愿祢可钦崇的荣耀受赞颂，从祢荣光的居所，基督，我们过犯和罪恶的平息者，祢藉着祢那著名、圣洁、赋予生命且神圣的奥迹，除去我们的愚妄。基督，我们本性的希望，从现在直到永远。阿们。\par

\Emph{封印或最后祝福。}\par

愿我们的主耶稣基督，我们曾事奉、瞻仰并尊崇的那一位，在祂那著名、圣洁、赋予生命且神圣的奥迹中，亲自使我们堪受祂王国辉煌的荣耀，并与祂的圣天使同乐，使我们在祂面前充满信心，得以站立在祂的右边。\par

但愿祂的仁慈与怜悯不断倾注于我们整个会众，从今时直到永远。\par

\Emph{主日和庆节。}\par

愿那藉着我们的主耶稣基督，以天上一切属神的祝福祝福了我们，为我们预备了祂的国度，召叫我们进入那永不停止、永不消逝的可羡慕的美好事物，正如祂在赋予生命的福音中所应许我们的，并对祂门徒的蒙福会众说：\BibleQuote{我实实在在告诉你们：谁吃我的肉，并喝我的血，便住在我内，我也住在他内，在末日我要使他复活；并且他不受审判，我要使他出死入生。}\par

愿祂现在亲自祝福这会众，保守我们的地位，荣耀我们的人民，他们前来并因领受祂那著名、圣洁、赋予生命且神圣的奥迹而喜乐；并愿你们因主十字架的圣号，从一切隐秘和公开的凶恶中，受封印并蒙保护，从今时直到永远。\par

\SourceRecord{Translated by James Donaldson.；Ante-Nicene Fathers；Vol. 7.；Edited by Alexander Roberts, James Donaldson, and A. Cleveland Coxe.；Buffalo, NY: Christian Literature Publishing Co.,；1886.；<http://www.newadvent.org/fathers/0719.htm>.}
